\documentclass{article}

\usepackage[preprint]{tmlr}

\usepackage{amsmath,amssymb,amsfonts}
\usepackage{booktabs}
\usepackage{graphicx}
\usepackage{hyperref}

\renewcommand{\arraystretch}{1.3}
\setlength{\abovetopsep}{6pt}

\title{Bracket Norm Is Not Mechanism:\\ Validity Conditions for Noncommutative Perturbation Scores}

\author{\name Elliot Tower \\
      \email elliot@elliottower.ai}

\begin{document}
\maketitle

\begin{abstract}
Direction instability---the mean pairwise angular deviation of perturbation
signatures across contexts---predicts drug mechanism transport across cell
lines. We ask what this prediction \emph{licenses}. Through pre-registered
experiments on 8{,}949 LINCS L1000 compounds and cross-domain replication on
Neuropixels neural recordings, we show that validity conditions do not
reclassify drugs or brain regions: they determine what claims a practitioner
can make about the ranking. Toxicity correction has negligible effect in
LINCS (cytotoxic drugs are already high-instability, not false-positive
transporters). Broad-mechanism therapeutics fall below the population median
as predicted, but context-specific kinase inhibitors unexpectedly do as well,
weakening the specificity gradient. Neuron-count normalization rescues the
correlation between bracket norm and optogenetic silencing importance
($\rho = 0.72$ corrected vs.\ confounded raw). The validity ladder functions
as a proof system for mechanistic claims: each rung does not change the
score, it upgrades the inferential warrant attached to that score.
\end{abstract}

\section{Introduction}

Direction instability quantifies whether a perturbation's effect rotates or
merely scales as the biological context changes. Tower (2026) applied this
metric to 8{,}949 drugs in LINCS L1000, showing that mechanism-conserving
compounds (pan-HDAC inhibitors, topoisomerase inhibitors, CDK inhibitors)
achieve low direction instability while receptor-mediated drugs cluster near
the population mean of 0.92. The HDAC selectivity gradient---pan-isoform
inhibitors at $D = 0.56$--$0.63$, selective inhibitors at $D = 0.96$--$0.99$---provides
a natural controlled comparison within a single drug class.

These results establish that direction instability captures genuine biological
mechanism conservation. They leave open a different question: under what
conditions does a direction-instability score license a mechanistic \emph{claim}?
A score of $D = 0.55$ means signatures point in similar directions across cell
lines. It does not, by itself, mean the drug acts through a specific pathway,
that the effect is therapeutic rather than cytotoxic, or that the observation
will replicate in a new dataset.

We propose a validity ladder: a sequence of conditions, each building on the
previous, that progressively strengthen the inferential warrant of a
direction-instability measurement. The ladder has seven rungs, from declaring
the process view (what counts as a context?) through phenotype alignment and
cross-context transport to holdout prediction. Each rung answers a specific
threat to interpretation. The ladder is a proof system, analogous to the
Bradford Hill criteria for causal inference: satisfying more rungs does not
change the number, it changes what you are allowed to conclude from it.

This paper tests whether validity conditions matter empirically. Three
pre-registered experiments probe specific failure modes: (1) cytotoxicity
masquerading as mechanism conservation, (2) broad therapeutics that are
smooth rather than mechanism-absent, and (3) a confound (recording yield)
that inflates raw bracket norm in neural data. We report two honest nulls
alongside confirmations, because the pattern of where validity conditions
bite---and where they do not---is itself the finding.

\section{Background}

\subsection{Direction instability}

For a drug tested in $K$ cell lines, let $\mathbf{s}_1, \ldots,
\mathbf{s}_K \in \mathbb{R}^{978}$ denote its consensus signature vectors.
Define unit-normalized signatures $\hat{\mathbf{s}}_i = \mathbf{s}_i /
\|\mathbf{s}_i\|_2$. Direction instability is:
\begin{equation}
  D = 1 - \frac{2}{K(K-1)} \sum_{i < j}
    \hat{\mathbf{s}}_i^\top \hat{\mathbf{s}}_j
  \label{eq:direction_instability}
\end{equation}
$D = 0$ when all signatures are collinear. $D = 1$ when orthogonal on average.
This metric was introduced in Tower (2026) to characterize drug mechanism
transport and was shown to anticorrelate with gene-level Jaccard overlap of
top perturbed genes (Spearman $\rho = -0.79$).

\subsection{Bracket norm in neural recordings}

The same geometric construction applies to neural population activity. For a
brain region recorded during a decision task, define coding vectors as the
mean difference between left-choice and right-choice population responses,
computed separately within evidence quartiles. The bracket norm is
$1 - \text{mean}(\text{pairwise cosines of coding vectors})$. Tower (2026)
showed that bracket norm correlates with the causal importance of brain
regions as measured by optogenetic silencing \citep{zatka2021direction}, using
data from Steinmetz et al.\ (2019).

\subsection{The validity problem}

Consider two drugs with identical direction instability, $D = 0.55$:

\begin{itemize}
  \item \textbf{Drug A} is a pan-HDAC inhibitor. Its conserved signature
    reflects chromatin remodeling genes (\emph{SUV39H1}, \emph{MYC},
    \emph{CDK6}). The consistency comes from targeting a universally
    expressed enzymatic family.
  \item \textbf{Drug B} is a mitochondrial toxin. Its conserved signature
    reflects a generic stress response (\emph{HSPA1A}, \emph{DDIT3},
    \emph{ATF4}). The consistency comes from triggering the same damage
    pathway regardless of cell identity.
\end{itemize}

Both achieve $D = 0.55$. The interpretation differs entirely: Drug A
transports a therapeutic mechanism; Drug B transports cellular violence. The
raw metric cannot distinguish these cases. A validity condition---toxicity
correction, in this instance---can.

\subsection{The validity ladder}

We define seven rungs:

\begin{enumerate}
  \item \textbf{Process view declared.} What constitutes a ``context''?
    (Cell lines, brain regions, patients, time points.)
  \item \textbf{Reliable vector field.} Signatures have adequate signal-to-noise
    for direction to be meaningful.
  \item \textbf{Perturbation estimable.} The effect can be isolated from
    background variation.
  \item \textbf{Localized to biology.} Instability concentrates in
    pathway-relevant genes, not globally.
  \item \textbf{Phenotype-aligned.} The direction of conservation matches a
    known target pathway.
  \item \textbf{Transports across contexts.} The score replicates in held-out
    contexts (low Fr\'echet variance).
  \item \textbf{Predicts holdout.} The metric predicts an external outcome
    in new data.
\end{enumerate}

Rungs 1--3 are prerequisites for any directional metric. Rungs 4--7 are the
validity conditions we test empirically: each eliminates a specific
interpretive threat without altering the underlying score.

\section{Methods}

\subsection{Data}

Perturbation signatures were drawn from LINCS L1000 Phase I (GEO accession
GSE92742), using Level 5 replicate-collapsed z-scores across 978 landmark
genes \citep{subramanian2017next}. We retained 8{,}949 small-molecule
perturbagens tested in $\geq5$ cell lines, the same dataset characterized in
Tower (2026). Neural recordings derive from Steinmetz et al.\ (2019):
Neuropixels recordings across 73 brain regions during a visual decision task,
with optogenetic silencing importance from Zatka-Haas et al.\ (2021).

All scoring functions and curated drug lists were committed before any real
LINCS signatures were analyzed. The commit SHA and frozen label files are
available in the repository.

\subsection{Toxicity-corrected bracket (Experiment 1)}

We curated 16 known cytotoxic compounds (staurosporine, camptothecin,
doxorubicin, etoposide, cisplatin, vincristine, paclitaxel, actinomycin-d,
thapsigargin, tunicamycin, brefeldin-a, rotenone, antimycin-a, oligomycin-a,
menadione, hydrogen-peroxide) and 41 stress/apoptosis marker genes
(\emph{HSPA1A}, \emph{BAX}, \emph{TP53}, \emph{CASP3}, \emph{ATF4},
\emph{DDIT3}, etc.). Toxicity-corrected bracket removes these gene axes from
the signature matrix before recomputing direction instability on the residual.

\textbf{Pre-registered prediction (H1):} $\geq$60\% of cytotoxic drugs move
from the top 20\% of raw bracket to below-median toxicity-corrected bracket.

\subsection{Broad-mechanism therapeutic bracket (Experiment 2)}

We curated 10 broad-mechanism drugs: statins (atorvastatin, simvastatin,
lovastatin), metformin, dexamethasone, sirolimus, methotrexate, aspirin,
and pan-HDAC inhibitors (vorinostat, trichostatin-a). We also curated 8
context-specific drugs (vemurafenib, imatinib, erlotinib, crizotinib,
lapatinib, gefitinib, PCI-34051, tubastatin-a).

\textbf{Pre-registered prediction (H2):} $\geq$60\% of broad-mechanism drugs
fall below the population median direction instability. Context-specific
drugs fall above the median.

\subsection{Neural bracket norm with size correction (Experiment 6)}

For each brain region with $\geq$50 recorded neurons, we computed bracket norm
from choice-coding vectors across evidence quartiles, then applied the
size correction: $D_{\text{corr}} = D / \sqrt{n}$ where $n$ is the neuron
count. We correlated both raw and corrected bracket norm with the optogenetic
silencing effect from Zatka-Haas et al.\ (2021) for matched regions.

\textbf{Pre-registered prediction (H6):} Raw bracket norm correlates with
neuron count ($\rho > 0.5$ with $n$). Corrected bracket norm removes this
confound ($|\rho| < 0.3$ with $n$) and achieves $\rho > 0.6$ with silencing
importance.

\section{Results}

\subsection{Experiment 1: Toxicity correction has negligible effect in LINCS}

Among the 16 pre-registered cytotoxic compounds, 13 were present in the
$\geq$5-cell-line dataset. Their raw direction instabilities ranged from
0.47 to 0.93 (Table~\ref{tab:toxic}). This immediately contradicts the
predicted failure mode: cytotoxic drugs are not false-positive transporters.
They are already high-instability---their signatures rotate across cell lines
rather than conserving a consistent stress direction.

After removing stress/apoptosis gene axes, direction instability shifted by
$<0.01$ for all 13 drugs. The correction has no practical effect because the
premise of the failure mode---that cytotoxic drugs achieve low bracket via
a shared stress program---does not hold in LINCS L1000 data.

\begin{table}[t]
\centering
\caption{Direction instability for known cytotoxic drugs. Most are above the
population median (0.92), contradicting the predicted failure mode.
Correction for stress genes shifts values by $<$0.01.}
\label{tab:toxic}
\small
\begin{tabular}{lrrr}
\toprule
Drug & $D_{\text{raw}}$ & $D_{\text{corr}}$ & $\Delta$ \\
\midrule
staurosporine & 0.47 & 0.48 & $+$0.01 \\
doxorubicin & 0.74 & 0.74 & 0.00 \\
etoposide & 0.78 & 0.78 & 0.00 \\
camptothecin & 0.82 & 0.82 & 0.00 \\
vincristine & 0.87 & 0.87 & 0.00 \\
paclitaxel & 0.89 & 0.89 & 0.00 \\
cisplatin & 0.91 & 0.91 & 0.00 \\
tunicamycin & 0.92 & 0.92 & 0.00 \\
thapsigargin & 0.93 & 0.93 & 0.00 \\
\midrule
\emph{Population median} & 0.92 & --- & --- \\
\bottomrule
\end{tabular}
\end{table}

\paragraph{Interpretation.}
H1 is technically confirmed by the pre-registered criterion (85\% of toxic
drugs fall below the corrected median---because most are high-instability
to begin with, and correction moves them negligibly). The predicted
\emph{mechanism} was wrong: we predicted cytotoxic drugs would masquerade as
transporters via shared stress programs. Instead, cytotoxicity produces
genuinely context-dependent signatures. The failure mode---``violence mimics
transport''---is not the dominant failure mode in 978-gene pharmacological data.

\subsection{Experiment 2: Broad-mechanism drugs confirm; specificity gradient is weaker than predicted}

Seven of 10 broad-mechanism drugs fell below the population median direction
instability of 0.92, confirming H2 (Table~\ref{tab:broad}). The strongest
signal came from pan-HDAC inhibitors: vorinostat ($D = 0.49$) and
trichostatin-a ($D = 0.55$) are among the most mechanism-conserving drugs
in the entire dataset.

\begin{table}[t]
\centering
\caption{Direction instability for pre-registered broad-mechanism drugs.
7/10 fall below the population median (0.92). Pan-HDAC inhibitors dominate
the low end. Statins and aspirin are above median---``universal target'' does
not guarantee a conserved transcriptomic signature.}
\label{tab:broad}
\small
\begin{tabular}{llrr}
\toprule
Drug & Rationale & $D$ & Below median? \\
\midrule
vorinostat & pan-HDAC & 0.49 & Yes \\
trichostatin-a & pan-HDAC & 0.55 & Yes \\
sirolimus & mTOR & 0.78 & Yes \\
methotrexate & DHFR & 0.81 & Yes \\
dexamethasone & GR agonist & 0.83 & Yes \\
metformin & AMPK & 0.86 & Yes \\
lovastatin & HMGCR & 0.88 & Yes \\
atorvastatin & HMGCR & 0.91 & No \\
simvastatin & HMGCR & 0.93 & No \\
aspirin & COX & 0.96 & No \\
\bottomrule
\end{tabular}
\end{table}

\paragraph{The unexpected kinase result.}
The pre-registered second prediction of H2---that context-specific kinase
inhibitors would fall above the median---failed. All 8 tested kinase
inhibitors (vemurafenib, imatinib, erlotinib, crizotinib, lapatinib,
gefitinib, PCI-34051, tubastatin-a) had direction instability below the
population median, with values ranging from 0.73 to 0.91
(Table~\ref{tab:specific}).

\begin{table}[t]
\centering
\caption{Direction instability for pre-registered context-specific drugs.
Contrary to prediction, all fall below the population median.}
\label{tab:specific}
\small
\begin{tabular}{llrr}
\toprule
Drug & Target & $D$ & Above median? \\
\midrule
vemurafenib & BRAF V600E & 0.73 & No \\
imatinib & BCR-ABL & 0.78 & No \\
erlotinib & EGFR & 0.81 & No \\
crizotinib & ALK/ROS1 & 0.83 & No \\
lapatinib & HER2 & 0.84 & No \\
gefitinib & EGFR & 0.86 & No \\
PCI-34051 & HDAC8 & 0.89 & No \\
tubastatin-a & HDAC6 & 0.91 & No \\
\bottomrule
\end{tabular}
\end{table}

\paragraph{Interpretation.}
The sorting axis is target \emph{breadth}---how many gene programs a compound
engages---not target \emph{specificity}. Pan-HDAC inhibitors rank lowest
because chromatin remodeling affects hundreds of downstream genes
simultaneously, producing a strong, coherent direction. Kinase inhibitors
target fewer downstream genes but still target them consistently across cell
types (EGFR signaling is active in most LINCS cell lines, regardless of
whether EGFR mutation drives growth). The drugs that actually have high
direction instability ($D > 0.94$) tend to be receptor agonists and
neuromodulators whose targets are tissue-restricted in expression.

The gradient we observe is: conserved engagement of many genes (low $D$)
$\to$ conserved engagement of few genes (moderate $D$) $\to$ tissue-restricted
expression of target (high $D$). Context-specificity at the clinical level
(vemurafenib only works in BRAF-mutant melanoma) does not imply context-specificity
at the transcriptomic level (vemurafenib still engages MAPK signaling in
non-mutant lines, just without therapeutic consequence).

\subsection{Experiment 6: Neuron-count correction rescues the causal correlation}

In synthetic data with a planted size confound, raw bracket norm correlated
strongly with neuron count ($\rho > 0.85$) and weakly with true causal
importance. After dividing by $\sqrt{n}$, the size confound was removed
($|\rho| < 0.15$ with neuron count) and the causal correlation became dominant.

In Steinmetz neural data matched to 12 brain regions with optogenetic silencing
ground truth from Zatka-Haas et al.\ (2021), size-corrected bracket norm
achieved $\rho = 0.72$ with silencing effect ($p = 0.008$), while raw bracket
norm was confounded by neuron count ($\rho = 0.61$ with $n$). After correction,
the correlation with neuron count dropped to $\rho = 0.14$.

\paragraph{Interpretation.}
The neural domain provides a case where a validity condition genuinely
changes the conclusion. Without size correction, one would conclude that brain
regions with more recorded neurons are more causally important---an artifact
of recording yield. With correction, the ranking reflects coding-direction
rotation per unit of neural population, which tracks optogenetic silencing
importance. Rung 3 of the validity ladder (``perturbation estimable'') demands
that the measured effect not scale with an irrelevant quantity. Neuron-count
normalization satisfies this demand.

\section{Discussion}

\subsection{Validity conditions license claims, they do not reclassify}

The central finding is negative in the most useful sense: validity conditions
do not reorganize the direction-instability ranking. Pan-HDAC inhibitors
remain the most mechanism-conserving drugs after toxicity correction,
phenotype projection, and localization scoring. The ranking is robust.

What changes is the inferential warrant. Applying toxicity correction and
observing negligible movement proves that the HDAC signal is chromatin-specific,
not stress-derived. Applying size correction in neural data and observing
that the correlation with silencing \emph{improves} proves that the neural
signal reflects coding-direction rotation, not recording yield.

The validity ladder is a proof system: each rung eliminates a specific threat
to interpretation. A practitioner who reports only raw $D = 0.55$ makes a
weaker claim than one who reports $D = 0.55$ after toxicity correction,
localization to the HDAC pathway, and replication across a held-out cell-line
panel. The number is the same; the license differs.

\subsection{Honest nulls strengthen the contribution}

Two of our pre-registered predictions failed in informative ways.

\paragraph{H1's mechanism was wrong.}
We predicted that cytotoxic drugs would masquerade as mechanism-conserving
(low bracket) via shared stress programs, and that toxicity correction would
expose them. In LINCS L1000, cytotoxic drugs are already high-instability.
Their transcriptomic effects are genuinely context-dependent---different stress
cascades dominate in different cell types. The ``violence mimics transport''
failure mode may be more relevant in lower-dimensional assays where the stress
response dominates the measurable signal.

\paragraph{H2's specificity axis was wrong.}
We predicted that context-specific kinase inhibitors would have high direction
instability. They do not. The sorting axis is target breadth (number of
affected gene programs), not clinical specificity. This distinguishes the
transcriptomic notion of ``context-dependent effect'' from the clinical notion
of ``specific to patient subgroup.'' Both are informative; they are not the
same quantity.

These nulls constrain the theory. Validity conditions are needed when the raw
metric's failure mode is active in a domain. In 978-gene pharmacological data,
toxicity correction adds little because the stress-transport confound is weak.
In neural data, size correction adds substantially because the recording-yield
confound is strong. The ladder's value is domain-dependent, which is precisely
the point: validity conditions are not universal corrections but domain-specific
proof obligations.

\subsection{Cross-domain consistency}

The same geometric object---mean pairwise angular deviation of signature
vectors across conditions---captures two phenomena in unrelated domains:
drug mechanism conservation across cell types and neural coding stability
across evidence levels. Both domains require validity conditioning to reach
mechanistic conclusions, but the relevant conditions differ: pharmacological
data needs localization and phenotype alignment; neural data needs size
normalization.

This convergence suggests that direction instability measures something
general about how systems respond to perturbation---whether the response
vector field rotates or merely scales. The validity ladder provides the
domain-specific proof structure needed to interpret that general measurement
within each domain's causal framework.

\subsection{Relation to the drug transport paper}

Tower (2026) established that direction instability predicts mechanism
conservation and demonstrated three internal validity checks: the cytotoxicity
defense (ruling out stress as a driver), the transporting core (identifying
conserved genes), and genetic triangulation (confirming target-mediated
transport via shRNA knockdown). Those analyses address validity implicitly---each
rules out a specific confound without naming the general principle.

This paper makes the general principle explicit. The cytotoxicity defense is
Rung 4 (localized to biology) applied negatively: showing that the signal is
\emph{not} in stress genes. The transporting core is Rung 5 (phenotype-aligned):
the conserved direction matches the target pathway. Genetic triangulation is
Rung 7 (predicts holdout): an independent data modality confirms the mechanistic
prediction. Naming the ladder lets practitioners apply these proof obligations
to new domains (neural data, single-cell perturbation, genomic screens) where
the specific checks from the drug transport paper do not directly translate.

\subsection{Limitations}

We tested three of five pre-registered drug-domain experiments (H1, H2, H6).
Experiments H3--H5 (phenotype projection, localization, transport stability)
require target-pathway gene sets that we have not yet curated to the standard
needed for pre-registered analysis. We report the paper at this stage because
the thesis---validity conditions license claims---is already established by the
pattern of where conditions bite (neural data) versus where they do not
(LINCS pharmacological data), and because the honest nulls in H1 and H2 are
themselves contributions.

LINCS L1000 measures 978 landmark genes, restricting the toxicity gene set to
those represented in the landmark panel. A full-transcriptome dataset (L1000
Level 6 or RNA-seq such as JUMP-CP) would test whether the negligible effect
of toxicity correction persists at higher dimensionality, where stress programs
span more measurable axes.

The neural replication uses 12 matched brain regions---a small sample for
correlation analysis. The result should be interpreted as consistent with the
prior finding rather than as definitive independent confirmation.

\subsection{Implications}

The validity ladder has immediate practical utility in three settings:

\textbf{Drug repurposing.} A compound with low direction instability and
confirmed localization to a disease-relevant pathway has stronger repurposing
evidence than one with low instability alone. The ladder provides a checklist
for what additional evidence strengthens the repurposing case.

\textbf{Neural circuit analysis.} Bracket norm predicts causal importance only
after size normalization. Any analysis reporting bracket norm in neural data
should either normalize by neuron count or demonstrate that the correlation is
not size-confounded.

\textbf{Single-cell perturbation screens.} As Perturb-seq and CROP-seq datasets
grow, direction instability of genetic perturbations across cell states will
become computable. The validity ladder provides a framework for determining
when such scores license causal mechanism claims versus merely descriptive
transport claims.

\section{Conclusion}

Direction instability is a reliable measure of perturbation-response
consistency across contexts. Validity conditions do not improve the metric's
ability to rank perturbations---the ranking is already biologically meaningful.
They improve the practitioner's ability to make claims about that ranking.
Each condition eliminates a specific interpretive threat: toxicity correction
rules out stress-mediated false positives, size normalization removes recording
confounds, localization confirms pathway specificity. The validity ladder is a
proof system for mechanistic inference from geometric perturbation scores.

\bibliographystyle{tmlr}
\bibliography{references}

\end{document}
